\section{Probability Space}

\subsection{Random Experiment}
\begin{itemize}
\item A \textbf{random experiment} is an experiment in which the
outcome varies in an unpredictable fashion when the experiment is
repeated under the same conditions, i.e., \emph{random}.
\item An \textbf{outcome} of a random experiment is a result that
cannot be decomposed into other results.
\end{itemize}

\subsection{Sample Space}
\begin{itemize}
\item The set of all possible outcomes for a random experiment is
called the \textbf{sample space} \(\Omega\).
\end{itemize}

\paragraph{Examples:}
\begin{enumerate}
\item Tossing a coin: \(\Omega=\{H, T\}\)
\item Tossing a coin twice: \(\Omega=\{H H, H T, T H, T T\}\)
\item Taking EEE 5544 and getting a grade:
\(
\Omega=\{A, A-, B+, B, B-, C+, C, D, F\}
\)
\item Picking an integer at random:
\(\Omega=\{\ldots,-3,-2,-1,0,1,2,3, \ldots \}\)
\item Checking the temperature (in $K$) at a random location
\(
\Omega=(0, \infty)
\)
\end{enumerate}

\begin{itemize}
\item \textbf{Finite} \(\Omega\): Examples 1, 2, 3
\item \textbf{Countably infinite} \(\Omega\): Example 4
\item \textbf{Discrete} \(\Omega\):  Examples 1--4
\item \textbf{Continuous (uncountably infinite)} \(\Omega\):  Example
5
\end{itemize}

\subsection{Event Space}
\begin{itemize}
\item An \textbf{event} is a set of outcomes (a subset of \(\Omega\))
   that we can ``quantify its happenstance'' (see below). 
\item Two special events: 
\begin{itemize}
\item \(\Omega =\)  \textbf{certain event} 
\item \(\emptyset =\) \textbf{impossible event}
\end{itemize}
\item An \textbf{event space} \(\mathcal{F}\) of a random experiment
is a collection of events.
\end{itemize}

\subsubsection{Discrete sample space}
\begin{itemize}
\item For a discrete \(\Omega\), we typically consider the event space
that is the collection
of all subsets of \(\Omega\), i.e., \(\mathcal{F} = 2^{\Omega}\),
which is called the \emph{power set} of \(\Omega\).
\end{itemize}

\paragraph{Examples:}
\begin{enumerate}
\item Tossing a coin: \(\mathcal{F} =\{\emptyset,\{H\},\{T\},\{H T\}\}\).
\item Tossing a coin twice:  
\begin{align*}
\mathcal{F} =& \left\{\emptyset,\{H H\},\{T T\},\{H T\},\{T H\},
 \{H H, T T\},\{H H, H T\}, \ldots, \right. \\
 & \left. \{H H, T T, H T\}, \cdots,\{H H, T T, H T, T H\} \right\}.
\end{align*}
\item Generally, if 
\( \Omega=\left\{a_{1}, a_{2}, \ldots
a_{N}\right\}\), then 
\begin{equation*}
\mathcal{F}=\left\{\emptyset,\{a_{1}\},
\ldots,\{a_{N}\}, \{a_{1}, a_{2}\},
\ldots,\{a_{1}, a_{2}, a_{3}\}, \ldots,
\{a_{1}, a_{2}, \ldots, a_{N}\}\right\}
\end{equation*}
Altogether \(2^{N}\) events (why?).
\end{enumerate}

  \item[-] What Roppen when \(\Omega\) is infuitely?\\
Countably infinite: - same idea applies \(F=2^{\Omega}\)
  \begin{itemize}
    \item[-] countably infinite events in 7.
    \item[-] need a better way \(T_{0}\) describe
  \end{itemize}
\end{itemize}

Uncountably infinite : Hum....?\\
Luckily rathematicians have solver the problem for us.

\begin{itemize}
  \item[-] Need a rath object to allow description of a random experiment.
  \item[-] Probability space \((\Omega, F, P)\) Sample space
  \item[] probability measure
  \item[-] Sample space \(\Omega\) : set that contains all possible ants of the random experiment.
\end{itemize}
